\chapter{Full Operational Semantics}\label{app:chap:full}
\section{Semantic Objects}\label{app:sec:dom}
\begin{center}
  \begin{bnf}[rccll]
    v : $\dom{Val}$ ::=
    | $k$ : constant values
    | $cl$ : closure
    | $C$ : component name
    | $cs$ : component spec
    | $[\Overline{s}]$ : view spec list
    | $\<\ell, p\>$ : setter closure
    ;;
    k : $\dom{Const}$ ::= $\<\>$ // $\TT$ // $\FF$ // $n$ : constant values
    ;;
    cl : $\dom{Clos}$ ::=
    | $\clos{x}{e}{\sigma}$ : closure
    ;;
    \sigma : $\dom{Env}$ ::= $[\Overline[\ell]{x \mapsto v}]$ : environment
    ;;
    cs : $\dom{ComSpec}$ ::=
    | $\<C, v\>$ : component spec
    ;;
    s : $\dom{ViewSpec}$ ::= $k$ // $cl$ // $[\Overline{s}]$ // $cs$ : view spec
    ;;
    \delta : $\dom{DefTable}$ ::= $[\Overline[\ell]{C \mapsto \deftabent{x}{e}}]$ : component def table
    ;;
  \end{bnf}
\end{center}
\begin{center}
  \begin{bnf}[rccll]
    \ell :in: $\mathbb{N}$ : label
    ;;
    p :in: $\dom{Path}$ : tree path
    ;;
    m : $\dom{TreeMem}$ ::= $[\Overline[\ell]{p \mapsto \pi}]$ : tree memory
    ;;
    \pi : $\dom{View}$ ::= $\view{cs}{\{\Overline{d}\}}{\rho}{q}{t}$ : state environment
    ;;
    d : $\dom{Decision}$ ::= \dec{Check} // \dec{Effect} : decision
    ;;
    t : $\dom{Tree}$ ::= $k$ // $cl$ // $[\Overline{t}]$ // $p$ : tree
    ;;
    \rho : $\dom{SttStore}$ ::= $[\Overline[\ell]{\ell \mapsto \sttstent{v}{q}}]$ : state store
    ;;
    q : $\dom{JobQ}$ ::= $[\Overline[\ell]{cl}]$ : job queue
    ;;
    \Sigma : $\dom{Context}$ ::= $m$ // $\pi$ : evaluation context
    ;;
    \phi : $\dom{Phase}$ ::= \phase{Init} // \phase{Succ} // \phase{Normal} : phase
    ;;
    \mu : $\dom{Mode}$ ::= $\rendermode$ \textrm{(rendered)} // $\checkmode$ \textrm{(check)} // $\eloopmode$  \textrm{(event loop)} : runtime mode
  \end{bnf}
\end{center}


\section{Operational Semantics}\label{app:sec:sem}
\begin{center}
  \hfill\fbox{$\<e, \delta\> \text{ or }\<t, m, \omega, \delta, \mu\> \smallstep \<t, m', \omega', \delta, \mu'\>$}
\begin{mathpar}
\inferrule[StepInit]{
  \Step<\phase{Normal}><->[[], []]{e}{s, [], \omega} \\
  [] \vdash \sem{init}(s) = \<t, m, \omega'\>
}{
  \<e, \delta\> \smallstep \<t, m, \omega \concat \omega', \delta, \rendermode\>
} \and
\inferrule[StepEffect]{
  m \vdash \sem{commitEffs}(t) = \<m', \omega'\>
}{
  \<t, m, \omega, \delta, \rendermode\> \smallstep \<t, m', \omega \concat \omega', \delta, \checkmode\>
} \and
\inferrule[StepCheck]{
  m, \delta \vdash \sem{check}(t) = \<\mu, m', \omega'\> \\
}{
  \<t, m, \omega, \delta, \checkmode\> \smallstep \<t, m', \omega \concat \omega', \delta, \mu\>
} \and
\inferrule[StepEvent]{
  \clos{x}{e}{\sigma} \in \sem{handlers}(m, t) \\
  \Step<\phase{Normal}><->[m, \sigma[x \mapsto \<\>]]{e}{v, m', \omega'} \\
}{
  \<t, m, \omega, \delta, \eloopmode\> \smallstep \<t, m', \omega \concat \omega', \delta, \checkmode\>
}
\end{mathpar}
\end{center}

\begin{center}
  % NOTE: handlers
  \hfill\fbox{$\sem{handlers}(m, t) = \{\Overline[\ell]{cl}\}$}
  \[
    \sem{handlers}(m, t) = \begin{dcases*}
      \{\} & if $t = k$ \\
      \{ cl \} & if $t = cl$ \\
      \bigcup_{i = 1}^n \sem{handlers}(m, t_i) & if $t = [\Overline{t_i}]_{i = 1}^n$ \\
      \sem{handlers}(m, m[p].\field{children}) & if $t = p$
    \end{dcases*}
  \]
\end{center}

\begin{center}
  \hfill\fbox{$\Step<\phi><p>[\Sigma, \sigma]{e}{v, \Sigma', \omega}$}
  \begin{mathparpagebreakable}
    \inferrule[Unit]{ }{\Step{\texttt{()}}{\<\>, \Sigma}, []} \and
    \inferrule[True]{ }{\Step{\texttt{true}}{\TT, \Sigma}, []} \and
    \inferrule[False]{ }{\Step{\texttt{false}}{\FF, \Sigma}, []} \and
    \inferrule[Int]{ }{\Step{n}{n, \Sigma}, []} \and
    \inferrule[Var]{ }{\Step{x}{\sigma(x), \Sigma}, []} \and
    \inferrule[Bop]{
      \Step[\Sigma, \sigma]{e_1}{n_1, \Sigma_1, \omega_1} \\
      \Step[\Sigma_1, \sigma]{e_2}{n_2, \Sigma_2, \omega_2} \\
      \text{$f_\oplus \in \mathbb{Z} \times \mathbb{Z} \to \mathbb{Z}$}
    }{\Step[\Sigma, \sigma]{e_1 \oplus e_2}{f_\oplus(n_1, n_2), \Sigma_2, \omega_1 \concat \omega_2}} \and
    \inferrule[Cond]{
      \Step[\Sigma, \sigma]{e_1}{b, \Sigma', \omega} \\
      \Step[\Sigma', \sigma]{\scond{b}{e_2}{e_3}}{v, \Sigma'', \omega'}
    }{
      \Step[\Sigma, \sigma]{\cond{e_1}{e_2}{e_3}}{v, \Sigma'', \omega \concat \omega'}
    } \and
    \inferrule[Func]{ }{\Step[\Sigma, \sigma]{\func{x}{e}}{\clos{x}{e}{\sigma}, \Sigma, []}} \and
    \inferrule[Seq]{
      \Step[\Sigma, \sigma]{e_1}{\_, \Sigma_1, \omega_1} \\
      \Step[\Sigma_1, \sigma]{e_2}{v_2, \Sigma_2, \omega_2}
    }{\Step[\Sigma, \sigma]{\seq{e_1}{e_2}}{v_2, \Sigma_2, \omega_1 \concat \omega_2}} \and
    \inferrule[List]{
      \bigl(\Step[\Sigma_{i-1}, \sigma]{e_i}{s_i, \Sigma_i, \omega_i}\bigr)_{i = 1}^n
    }{
      \Step[\Sigma_0, \sigma]{
        [\Overline{e_i}]_{i = 1}^n}{[\Overline{s_i}]_{i = 1}^n, \Sigma_n, {\textstyle\Concat_{i = 1}^n \omega_i}
      }
    } \and
    \inferrule[LetBind]{
      \Step{e_1}{v_1, \Sigma_1, \omega_1} \\
      \Step[\Sigma_1, \sigma[x \mapsto v_1]]{e_2}{v_2, \Sigma_2, \omega_2}
    }{
      \Step{\letbind{x}{e_1}{e_2}}{v_2, \Sigma_2, \omega_1 \concat \omega_2}
    } \and
    \inferrule[AppFunc]{
      \Step{e_1}{\clos{x}{e}{\sigma'}, \Sigma_1, \omega_1} \\
      \Step[\Sigma_1, \sigma]{e_2}{v_2, \Sigma_2, \omega_2} \\
      \Step[\Sigma_2, \sigma'[x \mapsto v_2]]{e}{v', \Sigma', \omega'}
    }{
      \Step{\app{e_1}{e_2}}{v', \Sigma', \omega_1 \concat \omega_2 \concat \omega'}
    } \and
    \inferrule[AppFunc]{
  \Step{e_1}{\clos{x}{e}{\sigma'}, \Sigma_1, \omega_1} \\
  \Step[\Sigma_1, \sigma]{e_2}{v_2, \Sigma_2, \omega_2} \\
  \Step[\Sigma_2, \sigma'[x \mapsto v_2]]{e}{v', \Sigma', \omega'}
}{
  \Step{\app{e_1}{e_2}}{v', \Sigma', \omega_1 \concat \omega_2 \concat \omega'}
} \and
\inferrule[AppCom]{
  \Step{e_1}{C, \Sigma_1, \omega_1} \\
  \Step[\Sigma_1, \sigma]{e_2}{v, \Sigma_2, \omega_2}
}{
  \Step{\app{e_1}{e_2}}{\<C, v\>, \Sigma_2, \omega_1 \concat \omega_2}
} \and
\inferrule[AppSetComp]{
  \Step[\pi, \sigma]{e_1}{\<\ell, p\>, \pi_1, \omega_1} \\
  \Step[\pi_1, \sigma]{e_2}{cl, \pi_2, \omega_2} \\
  \phi \in \{\phase{Init}, \phase{Succ}\}
}{
  \Step[\pi, \sigma]{\app{e_1}{e_2}}{
    \<\>,
    \pi_2\!\left\{
      \begin{spreadlines}{0pt}
        \begin{lgathered}
          \text{let $\rho = \pi_2.\field{sttst},\ vq = \rho[\ell]$ in} \\
          \begin{aligned}
            \field{dec} &\colon \pi_2.\field{dec} \cup \{\dec{Check}\} \\
            \field{sttst} &\colon \rho[\ell \mapsto vq\{\fv{sttq}{vq.\field{sttq} \concat [cl]}\}]
          \end{aligned}
        \end{lgathered}
      \end{spreadlines}
    \right\},
    \omega_1 \concat \omega_2
  }
} \and
\inferrule[AppSetNormal]{
  \Step<\phase{Normal}><->[m, \sigma]{e_1}{\<\ell, p\>, m_1, \omega_1} \\
  \Step<\phase{Normal}><->[m_1, \sigma]{e_2}{cl, m_2, \omega_2} \\
}{
  \Step<\phase{Normal}><->[m, \sigma]{\app{e_1}{e_2}}{
    \<\>,
    m_2\!\left[
      \begin{spreadlines}{0pt}
        \begin{lgathered}
          \text{let $\pi = m_2[p],\ \rho = \pi.\field{sttst},\ vq = \rho[\ell]$ in} \\
          p \mapsto \pi\!\left\{
            \begin{aligned}
              \field{dec} &\colon \pi.\field{dec} \cup \{\dec{Check}\} \\
              \field{sttst} &\colon \rho[\ell \mapsto vq\{\fv{sttq}{vq.\field{sttq} \concat [cl]}\}]
            \end{aligned}
          \right\}
        \end{lgathered}
      \end{spreadlines}
    \right],
    \omega_1 \concat \omega_2
  }
} \and
\inferrule[SttBind]{
  \Step<\phase{Init}>[\pi, \sigma]{e_1}{v_1, \pi_1, \omega_1} \\
  \Step<\phase{Init}>[{
    \begin{spreadlines}{0pt}
      \pi_1\!\left\{
        \fv{sttst}{
          \pi_1.\field{sttst}\!\left[\ell \mapsto \sttstent*{v_1}{[]}\right]
        }
      \right\}\!,
      \sigma\!\left[
        \begin{aligned}
          x &\mapsto v_1 \\
          x_{\textsf{set}} &\mapsto \<\ell, p\>
        \end{aligned}
      \right]
    \end{spreadlines}
  }]{e_2}{v_2, \pi_2, \omega_2}
}{
  \Step<\phase{Init}>[\pi, \sigma]{\stbind[\ell]{x}{e_1}{e_2}}{v_2, \pi_2, \omega_1 \concat \omega_2}
} \and
\inferrule[SttReBind]{
  \pi_0.\field{sttst}[\ell] = \bigl\{
    \fv{val}{v_0},\:\fv{sttq}{[\Overline{\clos{x_i}{e'_i}{\sigma_i}}]_{i = 1}^n}
  \bigr\} \\
  \bigl(\Step[\pi_{i-1}, \sigma_i[x_i \mapsto v_{i-1}]]{e'_i}{v_i, \pi_i, \omega_i}\bigr)_{i = 1}^n \\
  \Step[{
    \begin{spreadlines}{0pt}
      \pi_n\!\left\{
        \begin{aligned}
          \field{dec} &\colon \pi_n.\field{dec} \cup \scond{v_n \not\equiv v_0 }{\{\dec{Effect}\}}{\emptyset} \\
          \field{sttst} &\colon \pi_n.\field{sttst}[\ell \mapsto \sttstent{v_n}{[]}]
        \end{aligned}
      \right\}\!,
      \sigma\!\left[
        \begin{aligned}
          x &\mapsto v_n \\
          x_{\textsf{set}} &\mapsto \<\ell, p\>
        \end{aligned}
      \right]
    \end{spreadlines}
  }]{e_2}{v, \pi', \omega'}
}{
  \Step<\phase{Succ}>[\pi_0, \sigma]{\stbind{x}{e_1}{e_2}}{v, \pi', ({\textstyle\Concat_{i = 0}^n \omega_i}) \concat \omega'}
} \and
\inferrule[Eff]{
  \phi \in \{\phase{Init}, \phase{Succ}\}
}{
  \Step[\pi, \sigma]{\eff{e}}{
    \begin{spreadlines}{0pt}
      \<\>,
      \pi\{ \fv{effq}{\pi.\field{effq} \concat [\clos{\_}{e}{\sigma}]} \},
      []
    \end{spreadlines}
  }
} \and
\inferrule[Print]{
  \Step{e}{v, \Sigma', \omega}
}{
  \Step{\print{e}}{\<\>, \Sigma', \omega \concat [v]}
} \and
    \inferrule[Print]{
      \Step{e}{v, \Sigma', \omega}
    }{
      \Step{\print{e}}{\<\>, \Sigma', \omega \concat [v]}
    }
  \end{mathparpagebreakable}
\end{center}

\clearpage
\begin{center}
  \hfill\fbox{$\Step*<\phi><p>[\pi, \sigma]{e}{s, \pi', \omega}$}
\begin{mathparpagebreakable}
  \inferrule[EvalOnce]{
    \Step[\pi\{\fv{dec}{\pi.\field{dec}\setminus\{\dec{Check}\}},\fv{effq}{[]}\}, \sigma]{e}{s, \pi', \omega} \\
    \dec{Check} \notin \pi'.\field{dec} \\
    \phi \in \{\phase{Init}, \phase{Succ}\}
  }{
    \Step*[\pi, \sigma]{e}{s, \pi', \omega}
  } \and
  \inferrule[EvalMult]{
    \Step[\pi\{\fv{dec}{\pi.\field{dec}\setminus\{\dec{Check}\}},\fv{effq}{[]}\}, \sigma]{e}{s, \pi', \omega} \\
    \dec{Check} \in \pi'.\field{dec} \\
    \Step*<\phase{Succ}>[\pi', \sigma]{e}{s', \pi'', \omega'} \\
    \phi \in \{\phase{Init}, \phase{Succ}\}
  }{
    \Step*[\pi, \sigma]{e}{s', \pi'', \omega \concat \omega'}
  }
\end{mathparpagebreakable}
\end{center}

\begin{center}
  \hfill\fbox{$m, \delta \vdash \sem{init}(s) = \<t, m', \omega\>$}
\begin{mathparpagebreakable}
  \inferrule[InitConst]{ }{m, \delta \vdash \sem{init}(k) = \<k, m, []\>} \and
  \inferrule[InitClos]{ }{m, \delta \vdash \sem{init}(cl) = \<cl, m, []\>} \and
  \inferrule[InitArray]{
    \bigl(m_{i-1}, \delta \vdash \sem{init}(s_i) = \<t_i, m_i, \omega_i\>\bigr)_{i = 1}^n
  }{
    m_0, \delta \vdash \sem{init}([\Overline{s_i}]_{i = 1}^n) = \<[\Overline{t_i}]_{i = 1}^n, m_n, {\textstyle\Concat_{i = 1}^n \omega_i}\>
  } \and
  \inferrule[InitCom]{
    m \vdash p\:\judge{fresh} \\
    \delta[C] = \deftabent{x}{e} \\
    \Step*<\phase{Init}>[
      \view{\<C, v\>}{\emptyset}{[]}{[]}{\<\>},
      [x \mapsto v]
    ]{e}{s, \pi, \omega} \\
    m[p \mapsto \pi], \delta \vdash \sem{init}(s) = \<t, m', \omega'\>
  }{
    m, \delta \vdash \sem{init}(\<C, v\>) = \bigl\<p, m'[p \mapsto \pi\{\fv{dec}{\{\dec{Effect}\}}, \fv{child}{t}\}], \omega \concat \omega'\bigr\>
  }
\end{mathparpagebreakable}
\end{center}

\clearpage
\begin{center}
  \hfill\fbox{$m, \delta \vdash \sem{check}(t) = \<\mu, m', \omega\>$}
\begin{mathparpagebreakable}
  \inferrule[CheckConst]{ }{m, \delta \vdash \sem{check}(k) = \<\eloopmode, m, []\>} \and
  \inferrule[CheckClos]{ }{m, \delta \vdash \sem{check}(cl) = \<\eloopmode, m, []\>} \and
  \inferrule[CheckArray]{
    \bigl(m_{i-1}, \delta \vdash \sem{check}(t_i) = \<\mu_i, m_i, \omega_i\>\bigr)_{i = 1}^n
  }{\textstyle
    m_0, \delta \vdash \sem{check}([\Overline{t_i}]_{i = 1}^n) = \<\bigsqcup_{i = 1}^n \mu_i, m_n, {\textstyle\Concat_{i = 1}^n \omega_i}\>
  } \and
  \inferrule[CheckIdle]{
    m[p] = \pi \\
    \dec{Check} \notin \pi.\field{dec} \\
    m, \delta \vdash \sem{check}(\pi.\field{child}) = \<\mu, m', \omega\>
  }{\textstyle
    m, \delta \vdash \sem{check}(p) = \<\mu, m', \omega\>
  } \and
  \inferrule[CheckNoEffect]{
    m[p] = \pi \\
    \dec{Check} \in \pi.\field{dec} \\
    \pi.\field{spec} = \<C, v\> \\
    \delta[C] = \deftabent{x}{e} \\
    \Step*<\phase{Succ}>[
      \pi,
      [x \mapsto v]
    ]{e}{\_, \pi', \omega} \\
    \dec{Effect} \notin \pi'.\field{dec} \\
    m, \delta \vdash \sem{check}(\pi.\field{child}) = \<\mu, m', \omega'\>
  }{\textstyle
    m, \delta \vdash \sem{check}(p) = \<\mu, m'[p \mapsto \pi'], \omega \concat \omega'\>
  } \and
  \inferrule[CheckEffect]{
    m[p] = \pi \\
    \dec{Check} \in \pi.\field{dec} \\
    \pi.\field{spec} = \<C, v\> \\
    \delta[C] = \deftabent{x}{e} \\
    \Step*<\phase{Succ}>[
      \pi,
      [x \mapsto v]
    ]{e}{s, \pi', \omega} \\
    \dec{Effect} \in \pi'.\field{dec} \\
    m, \delta \vdash \sem{reconcile}(\pi.\field{child}, s) = \<t', m', \omega'\> \\
  }{\textstyle
    m, \delta \vdash \sem{check}(p) = \<\rendermode, m'[p \mapsto \pi'\{\fv{child}{t'}\}], \omega \concat \omega'\>
  }
\end{mathparpagebreakable}
\end{center}

\begin{center}
  \hfill\fbox{$m, \delta \vdash \sem{reconcile}(t, s) = \<t', m', \omega\>$}
\begin{mathparpagebreakable}
  % TODO: Update the rule
  \inferrule[ReconcileArray]{
    \bigl(m_{i-1}, \delta \vdash \sem{reconcile}(t_i, s_i) = \<t_i', m_i, \omega_i\>\bigr)_{i = 1}^n
  }{
    m_0, \delta
    \vdash \sem{reconcile}([\Overline{t_i}]_{i = 1}^n, [\Overline{s_i}]_{i = 1}^n)
    = \<[\Overline{t_i'}]_{i = 1}^n, m_n, {\textstyle\Concat_{i = 1}^n \omega_i}\>
  } \and
  \inferrule[ReconcileComEffect]{
    m[p] = \pi \\
    \pi.\field{spec} = \<C, \_\> \\
    \delta[C] = \deftabent{x}{e} \\
    \Step*<\phase{Succ}>[
      \pi\{\fv{spec}{\<C, v\>}\},
      [x \mapsto v]
    ]{e}{s, \pi', \omega} \\
    m, \delta \vdash \sem{reconcile}(\pi.\field{child}, s) = \<t', m', \omega'\> \\
  }{
    m, \delta \vdash \sem{reconcile}(p, \<C, v\>) = \bigl\<p, m'[p \mapsto \pi'\{\fv{dec}{\{\dec{Effect}\}}, \fv{child}{t'}\}], \omega \concat \omega'\bigr\>
  } \and
  \inferrule[ReconcileComNew]{
    m[p].\field{spec} = \<C', \_\> \\
    C \ne C' \\
    m, \delta \vdash \sem{init}(\<C, v\>) = \<t', m', \omega\>
  }{
    m, \delta \vdash \sem{reconcile}(p, \<C, v\>) = \<t', m', \omega\>
  } \and
  \inferrule[ReconcileOther]{
    \<t, s\> \neq \<[\Overline{t}], [\Overline{s}]\> \\
    \<t, s\> \neq \<p, \<C, v\>\> \\
    m, \delta \vdash \sem{init}(s) = \<t', m', \omega\>
  }{
    m, \delta \vdash \sem{reconcile}(t, s) = \<t', m', \omega\>
  }
\end{mathparpagebreakable}
\end{center}

\begin{center}
  \hfill\fbox{$m \vdash \sem{commitEffs}(t) = \<m', \omega\>$}
\begin{mathparpagebreakable}
  \inferrule[CommitEffsConst]{ }{m \vdash \sem{commitEffs}(k) = \<m, []\>} \and
  \inferrule[CommitEffsClos]{ }{m \vdash \sem{commitEffs}(cl) = \<m, []\>} \and
  \inferrule[CommitEffsArray]{
    \bigl(m_{i-1} \vdash \sem{commitEffs}(t_i) = m_i, \omega_i\bigr)_{i = 1}^n
  }{
    m_0 \vdash \sem{commitEffs}([\Overline{t_i}]_{i = 1}^n) = \<m_n, {\textstyle\Concat_{i = 1}^n \omega_i}\>
  } \and
  \inferrule[CommitEffsPathIdle]{
    \dec{Effect} \notin m[p].\field{dec} \\
    m \vdash \sem{commitEffs}(m[p].\field{child}) = \<m', \omega\> \\
  }{
    m \vdash \sem{commitEffs}(p) = \<m', \omega\>
  } \and
  \inferrule[CommitEffsPath]{
    \dec{Effect} \in m[p].\field{dec} \\
    m[p].\field{effq} = [\Overline{\clos{\_}{e_i}{\sigma_i}}]_{i = 1}^n \\
    m \vdash \sem{commitEffs}(m[p].\field{child}) = \<m_0, \omega_0\> \\
    \bigl(
      \Step<\phase{Normal}><->[m_{i-1}, \sigma_i]{e_i}{v_i, m_i, \omega_i}
    \bigr)_{i = 1}^n
  }{
    m \vdash \sem{commitEffs}(p) = \bigl\<m_n\bigl[p \mapsto m_n[p]\{\fv{dec}{m_n[p].\field{dec} \setminus \{\dec{Effect}\}}\}\bigr], {\textstyle\Concat_{i = 0}^n \omega_i}\bigr\>
  }
\end{mathparpagebreakable}
\end{center}


\chapter{Proofs of Theorems}\label{app:chap:proofs}
\begin{defn}[$e$-Equivalent Views]\label{def:equiv}
  Views~$\pi$ and~$\pi'$ are $e$-equivalent when the state entries whose labels appear in~$e$ are the same.
  That is,
  \[
    \pi \equiv_e \pi'
    \triangleiff
    \forall \ell \in \sem{labels}(e),\ \pi.\field{sttst}[\ell] = \pi'.\field{sttst}[\ell]. \qedhere
  \]
\end{defn}

\begin{defn}[$t$-Equivalent Memories]\label{def:memequiv}
  Tree memories~$m$ and~$m'$ are $t$-equivalent iff the descendant views of~$t$ in~$m$ and~$m'$ are the same.
  That is,
  \[
    m \equiv_t m'
    \triangleiff
    \forall p \in \sem{reachable}(m,t), m[p] = m'[p] \qedhere
  \]
\end{defn}

\begin{lem}[Similar Evaluations]\label{lem:simeval}
  Evaluations of~$e$ in~\phase{Succ} phase with similar views as a context produce identical values and output buffers, along with $e$-equivalent views.
  That is, if $\Step<\phase{Succ}>[\pi,\sigma]{e}{v,\pi',\omega}$,
  then for any $\hat\pi \approx \pi$, we have
  $\Step<\phase{Succ}>[\hat \pi,\sigma]{e}{v,\hat \pi',\omega}$ where $\hat\pi' \equiv_e \pi'$.
\end{lem}

\begin{proof}
  Suppose we have $\Step<\phase{Succ}>[\pi,\sigma]{e}{v,\pi',\omega}$ and choose some $\hat\pi \approx \pi$.
  We need to show that $\hat v = v$, $\hat\pi' \equiv_e \pi'$, and $\hat\omega = \omega$, where $\Step<\phase{Succ}>[\hat\pi,\sigma]{e}{\hat v,\hat\pi',\hat\omega}$.

  We prove that $\hat v = v$, $\hat\pi' \approx \pi'$ (before showing $\hat\pi' \equiv_e \pi'$), and $\hat\omega = \omega$ by rule induction.
  We put a~$\hat\cdot$ on the intermediate variables in the derivation of~$\Step<\phase{Succ}>[\hat\pi,\sigma]{e}{\hat v,\hat\pi',\hat\omega}$.
  \begin{description}
    \item[{\normalfont\Rule{Print}}]
      Evaluations of~$e$ produce same values~$\hat v = v$ and output buffers~$\hat \omega = \omega$ with similar views~$\hat\pi' \approx \pi'$ by the IH.
      Appending the same values to the same output buffers results in the same buffers~$\hat\omega \concat [\hat v] = \omega \concat [v]$.
      Resulting values are both~$\<\>$.
      Contexts---views in this case---are unmodified.

    \item[{\normalfont\Rule{AppSetComp}}]
      Evaluations of~$e_1$ and~$e_2$ produce same values~$\<\hat\ell, \hat p\> = \<\ell, p\>$ and $\hat{cl} = cl$, and output buffers~$\hat\omega_1 = \omega_1$ and $\hat\omega_2 = \omega_2$, and similar views~$\hat\pi_1 \approx \pi_1$ and $\hat\pi_2 \approx \pi_2$ by the IH.
      Concatenations of the same output buffers result in the same buffers~$\hat\omega_1 \concat \hat\omega_2 = \omega_1 \concat \omega_2$.
      Appending the same closures~$\hat{cl} = cl$ to the state update queues of similar views~$\hat\pi_2 \approx \pi_2$ produces similar views.
      (The added \dec{Check}s are discarded anyway during normalization.)
      Resulting values are both~$\<\>$.

    \item[{\normalfont\Rule{SttReBind}}]
      Evaluations of the queued~$e_i'$s all produce the intermediate values~$\hat v_i = v_i$ and output buffers~$\hat \omega_i = \omega_i$ that are equivalent and views~$\hat\pi_i \approx \pi_i$ that are similar by the IH.
      Then the modified~$\hat\pi_n$ and~$\pi_n$ each modified with the added \dec{Effect} and the same state~$\hat v_n = v_n$ are still similar.
      This suffices to apply the IH to the evaluation of~$e_2$, resulting in the same values~$\hat v = v$ and output buffers~$\hat\omega' = \omega'$, and similar views~$\hat\pi' \approx \pi'$.
      Thus the resulting values~$\hat v = v$ and the concatenation of the buffers~$(\concat_{i=1}^n \hat \omega_i) \concat \hat\omega' = (\concat_{i=1}^n \omega_i) \concat \omega'$ are equivalent, and the resulting views~$\hat\pi' \approx \pi'$ are similar.
  \end{description}
  The remaining cases are trivial.

  It is easy to show that if~$\pi$ and~$\hat \pi$ are similar and $e$-equivalent for some~$e$, then the results are also $e$-equivalent using rule induction.

  Now we can prove that $\hat\pi' \equiv_e \pi'$, again using rule induction.
  In each case, evaluation of each sub-expression~$e'$ of~$e$ produces $e'$-equivalent views by the IH, and the equivalence is preserved in following evaluations of other sub-expressions.
  Therefore, the resulting views from~$\pi$ and~$\hat \pi$ are $e'$-equivalent for every sub-expression~$e'$, which makes the resulting~$\pi'$ and~$\hat \pi'$ $e$-equivalent.
  Note that nested or dead sub-expressions of~$e$ do not contain state labels~$\ell$, due to the syntactic restriction that Hooks can only appear at the top level of a component body.

  The only noteworthy case is \Rule{SttReBind} where a state label exists.
  In this case, the state store entry for~$\ell$ is updated directly to the same value~$v_n$ (as just shown above) in both derivations from~$\pi$ and~$\hat \pi$.
  Therefore, the resulting~$\pi'$ and~$\hat \pi'$ have the same state store entry for~$\ell$.
\end{proof}

\simevallem*

\begin{proof}
  Let $\Step*<\phase{Succ}>[\pi,\sigma[x\mapsto v]]{e}{v,\pi',\omega}$ and $\Step*<\phase{Succ}>[\hat \pi,\sigma[x\mapsto v]]{e}{\hat v',\hat \pi', \hat \omega}$, for some~$\pi$ such that~$\pi.\field{spec} = \<C, v\>$ where~$\delta[C] = \<\lfun{x}{e}, \sigma\>$ and~$\hat\pi \approx \pi$.
  By \cref{lem:simeval}, evaluations of~$e$ produce $e$-equivalent views, that is, the states whose labels appear in~$e$ are the same in~$\pi'$ and~$\hat \pi'$.
  Since~$\pi'$ and~$\hat \pi'$ are valid under~$\delta$, the only labels that~$\pi'$ and~$\hat \pi'$ have are those appearing in~$e$.
  Therefore, $\pi' = \hat \pi'$.
  Similarly, $v' = \hat v'$ and $\omega = \hat \omega$.
  Any remaining loops, if present, are evaluated under identical views, trivially producing the same results.
\end{proof}

\begin{defn}[Stability and Semi-Stability]\label{def:stability}
  A view~$\pi$ is \emph{stable} iff the re-evaluation of~$\pi$'s body (as in \Rule{EvalOnce} and \Rule{EvalMult}) produces the same view as $\pi$:
  \begin{multline*}
    \Step[\pi\{\fv{dec}{\pi.\field{dec}\setminus\{\dec{Check}\}},\fv{effq}{[]}\}, \sigma[x\mapsto v]]{e}{\_, \pi, \omega} \\
    \text{where $\pi.\field{spec} = \<C,v\>$ and $\delta[C] = \<\lambda x.e,\sigma\>$.}
  \end{multline*}

  View~$\pi$ is \emph{semi-stable} iff $\pi$ with empty state update queues is stable.

  We extend this notion to tree memory:
  $m$ is~\emph{$t$-stable} iff the descendant views of~$t$ in~$m$ are all stable, and \emph{$t$-semi-stable} iff the descendant views of~$t$ in~$m$ are all semi-stable.

  Since component body of a stable view can be re-evaluated in an idempotent way, evaluation of the body can be thought of as ``reading'' the specification of the view.
\end{defn}

\begin{defn}[Coherence]\label{def:coherence}
  A view is \emph{coherent} if its decision is coherent with the status of the state update queue.
  That is, a coherent view~$\pi$ satisfies the predicate
  \[
    \dec{Check} \in \pi.\field{dec}
    \qquad\longleftrightarrow\qquad
    \exists \ell \in \domain \rho,\ \rho[\ell].\field{sttq} \neq []
    \quad \text{where $\rho = \pi.\field{sttst}$.}
  \]

  We extend this notion to tree memory, i.e., $m$ is \emph{$t$-coherent} iff the descendants of~$t$ in~$m$ are coherent.
\end{defn}

\begin{restatable}[Preservation of Validity]{lem}{validprop}\label{lem:valid}
  If $\<e,\delta\> \smallstep^* \<t,m,\omega,\delta,\mu\>$, then~$m$ is valid under~$\delta$.
\end{restatable}

\begin{proof}
  The only rule that adds states is \Rule{SttBind},
  which appears only in the derivation of $\sem{init}$ where the context being a fresh view.
  Therefore, it is sufficient to check that the evaluation of the component body in $\phase{Init}$ produces a valid view.
  Due to the syntactic restriction on the usage of Hooks, every |useState| Hook is executed during the evaluation, adding its state labeled~$\ell$ to the view.
  This process adds every label in the expression to the view, producing a valid view.
\end{proof}

\begin{lem}[Semi-Stability and Normalized Form]\label{lem:stabnorm}
  Let~$\pi$ be a semi-stable view valid under~$\delta$, and let~$\hat\pi = \sem{normalize}(\pi)$.
  If $\dec{Check} \notin \hat\pi.\field{dec}$, then the evaluation of the body produces a view exactly the same as the normalized form.
  That is, \[\Step*<\phase{Succ}>[\pi\{\fv{dec}{\pi.\field{dec} \setminus \{\dec{Check}\}},\fv{effq}{[]}\}, \sigma[x\mapsto v]]{e}{\_,\hat\pi, \omega}\] where~$\pi.\field{spec} = \<C,v\>$ and~$\delta[C] = \<\lfun x e, \sigma\>$.
\end{lem}

\begin{proof}
  The normalized view~$\hat\pi$ is the same as~$\pi$ except that the decision and state update queues are empty---$\dec{Check} \notin \hat\pi.\field{dec}$ implies also $\dec{Effect} \notin \hat\pi.\field{dec}$ from \cref{def:normalization}.
  Thus~$\hat\pi$ is stable as~$\pi$ is semi-stable.
  By \cref{def:stability}, evaluation with~$\hat\pi$ produces~$\hat\pi$ again, which also holds for retrying evaluation.
  Since~$\pi \approx \hat\pi$, by \cref{lem:simevalbody}, $\pi' = \hat\pi$.
\end{proof}


\begin{lem}[Stability of Retrying Evaluation]\label{lem:stab}
  Let~$\pi.\field{spec} = \<C,v\>$ and~$\delta[C] = \<\lambda x.e,\sigma\>$.
  If~$\pi$ is semi-stable and $\Step*[\pi,\sigma[x\mapsto v]]{e}{v',\pi', \omega}$,
  then~$\pi'$ is stable.
\end{lem}

\begin{proof}
  It is sufficient to show that the last evaluation of the component body produces a stable view.
  Since~$\dec{Check} \notin \pi'.\field{dec}$, \Rule{AppSetComp} is not included in the derivation and all state updates in~$\pi$ are pure.
  This means~$\sem{normalize}(\pi)$ has empty queues.
  Since~$\pi \approx \sem{normalize}(\pi)$, by the \cref{lem:stabnorm}, $\pi' = \sem{normalize}(\pi)$,
  which is stable.
\end{proof}

\begin{lem}[Preservation of Semi-Stability]\label{lem:semistab}
  If $\<e,\delta\> \smallstep^* \<t,m,\omega,\delta,\mu\>$, then~$m$ is $t$-semi-stable.
\end{lem}

\begin{proof}
  It is sufficient to show that the initial step generates a semi-stable tree memory, and derivation of each step preserves the semi-stability.
  \begin{description}
    \item[{\normalfont\Rule{StepInit}, \Rule{StepCheck}}]
      By \cref{lem:stab}, evaluation of each component body produces stable view.
      Updating views' children or adding decisions does not break the stability, so views modified by $\sem{init}$ and $\sem{visit}$ are always stable.
    \item[{\normalfont\Rule{StepEffect}, \Rule{StepEvent}}]
      The only rule modifying the view during normal evaluation is \Rule{AppSetNormal}.
      It only adds state updates or modify the decision of a view, which preserves semi-stability of the view. \qedhere
  \end{description}
\end{proof}

\begin{remark}
  While the semi-stability of memory is always preserved, the stability is often violated.
  When setter closures are applied during the execution of Effects or event handlers,
  the update functions are queued,
  which are flushed and applied during the next check.
  Note that the memory is indeed always stable after the check.
\end{remark}

\begin{restatable}[Preservation of Coherence]{lem}{coherence}\label{lem:coherence}
  If $\<e,\delta\> \smallstep^* \<t,m,\omega,\delta,\mu\>$, then~$m$ is $t$-coherent.
  That is, $\sem{init}$ and $\sem{visit}$ always produce~$t$-coherent tree memory.
\end{restatable}

\begin{proof}
  It is sufficient to show that the initial step generates a coherent tree memory, and the derivation of each step preserves the coherence.
  \begin{description}
    \item[{\normalfont\Rule{StepInit}}]
      A retrying evaluation of the component body produces a view whose decision does not include $\dec{Check}$ and state queues are empty, thus coherent.
      The result of $\sem{init}$ only contains new views.
      Therefore, the resulting memory $m$ is $t$-coherent.
    \item[{\normalfont\Rule{StepCheck}}]
      As stated above, the views whose bodies are evaluated are coherent.
      Views that are not evaluated are left as is in $m'$, thus preserving coherence.
      Therefore, each coherence of the views in the resulting memory $m'$ is preserved.
    \item[{\normalfont\Rule{StepEffect}, \Rule{StepEvent}}]
      The only rule modifying the view during normal evaluation is \Rule{AppSetNormal}. It adds a state update and adds $\dec{Check}$ decision to the view, which preserves the coherence of the view. \qedhere
  \end{description}
\end{proof}



\begin{lem}[Similar Reconciliations]\label{lem:reconsim}
  Reconciling with similar memories produces the equivalent results.
  That is, if $m \vdash \sem{reconcile}(t,s)=\<t',m', \omega\>$,
  then for $\hat m$ such that $m \approx_t \hat m$,
  $\hat m \vdash \sem{reconcile}(t,s) = \<t',\hat m', \omega\>$ and
  $m' \equiv_{t'} \hat m'$.
\end{lem}

\begin{proof}
  We proceed with induction on the derivation of $m \vdash \sem{reconcile}(t,s)=\<t',m', \omega\>$.
  \begin{description}
    \item[{\normalfont\Rule{ReconcileComNew}, \Rule{ReconcileOther}}]
      \sem{init}(s) does not read or modify existing views in $m$; it only adds new views to $m$.
      All new views are descendants of $t'$, and are the same in $m'$ and $\hat m'$.
      Since the evaluation contexts are identical, the output buffers are also the same.
    \item[{\normalfont\Rule{ReconcileList}}]
      Each derivation of $\sem{reconcile}(t_i,s_i)$ produces the same tree, the same output buffer, and $t'_i$-equivalent memories.
      Therefore, the final trees and the output buffers are the same and the final memories are $[\Overline{t'_i}]_{i=1}^n$-equivalent.
    \item[{\normalfont\Rule{ReconcileComEffect}}]
      The derivation of $\Bigstep*{\phi}{p}$ produces the same $v$, $\pi'$, and $\omega$ by \cref{lem:simevalbody} and $\sem{reconcile}(t,s)$ produces the same tree and output buffer with $t'$-equivalent memories by IH.
      Therefore, the final memories are also $t'$-equivalent. \qedhere
  \end{description}
\end{proof}

\begin{lem}[Similar Checks]\label{lem:simvisit}
  Checking with similar memories produces the equivalent results.
  That is, for similar memories $m$ and $\hat m$ that are $t$-coherent and $t$-semi-stable and have no view with \dec{Effect} decision set,
  if $m, \delta \vdash \sem{check}(t)=\<\mu,m', \omega\>$,
  then $\hat m, \delta \vdash \sem{check}(t) = \<\mu,\hat m', \omega'\>$ and
  $m' \equiv_{t'} \hat m'$.
  We also have $\omega = \omega'$ if component bodies do not print.
\end{lem}

\begin{proof}
  We proceed with induction on the derivation of $m, \delta \vdash \sem{check}(t)=\<\mu,m', \omega\>$.
  \begin{description}
    \item[{\normalfont\Rule{CheckConst}, \Rule{CheckClos}}] These cases are trivial.
    \item[{\normalfont\Rule{CheckList}}] ($t = [\Overline{t_i}]_{i=1}^n$).
      Each derivation of $\sem{check}(\hat t_i)$ produces memory, $t'_k$-equivalent to~$m_i$ for~$1 \le k \le i$ and $t'_k$-similar to $m_i$ for $i < k \le n$ by the IHs.
      As a result, $m_n$ and $\hat m_n$ are $t'_k$-equivalent for~$1 \le k \le n$, therefore $[\Overline{t'_i}]_{i=1}^n$-equivalent.
    \item[{\normalfont\Rule{CheckIdle}}] ($t = p$, $m[p] = \pi$, and $\hat m[p] = \hat \pi$).
      Note that $\pi$ is coherent and $\dec{Check} \notin \pi.\field{dec}$, implying that the state update queues are empty.
      Hence, $\sem{normalize}(\pi) = \pi$.
      \begin{itemize}
        \item $\dec{Check} \notin \hat \pi.\field{dec}$:
          By the IH, $\sem{check}(t)$ produces $t$-equivalent results.
        \item $\dec{Check} \in \hat \pi.\field{dec}$:
          By \cref{lem:stabnorm}, $\hat \pi' = \sem{normalize}(\hat \pi)$.
          Since $\pi \approx \hat \pi$, we have $\sem{normalize}(\hat \pi) = \sem{normalize}(\pi)$, and thus~$\hat\pi' = \pi$.
          Since no view, including~$\pi$, has \dec{Effect} in~$m$, $\dec{Effect} \notin \hat \pi'.\field{dec}$ as well.
          Therefore, $\hat m, \delta \vdash \sem{check}(p) = \_$ is derived from \Rule{CheckNoEffect},
          and thus $\hat m, \delta \vdash \sem{check}(t) = \<b, \hat m', \omega\>$.
          By the IH, $m' \equiv_t \hat m'$.
          As a result, the results are $p$-equivalent.
          Note that we have $\omega = \omega'$ if the component prints nothing.
      \end{itemize}
    \item[{\normalfont\Rule{CheckNoEffect}}] ($t = p$, $m[p] = \pi$, and $\hat m[p] = \hat \pi$).
      \begin{itemize}
        \item $\dec{Check} \notin \hat \pi.\field{dec}$:
          The proof is similar to the case of \Rule{CheckIdle} where $\dec{Check} \in \hat m[p].\field{dec}$, except that $m$ and $\hat m$ are swapped.
        \item $\dec{Check} \in \hat \pi.\field{dec}$:
          By \cref{lem:simevalbody}, $\hat \pi' = \pi'$ (and $\omega = \omega'$ if componets print nothing),
          and by the IH, $m' \equiv_p \hat m'$.
          Therefore, the results are $p$-equivalent.
      \end{itemize}
    \item[{\normalfont\Rule{CheckEffect}}] ($t = p$).
      We have $\dec{Check} \in \hat \pi.\field{dec}$ since $\sem{normalize}(\hat \pi) = \sem{normalize}(\pi)$ and $\dec{Check} \in \sem{normalize}(\pi).\field{dec}$.
      By \cref{lem:simevalbody}, $\pi' = \hat \pi'$,
      and by \cref{lem:reconsim}, $m' \equiv_p \hat m'$.
      Therefore, the results are $p$-equivalent. \qedhere
  \end{description}
\end{proof}

\simtransthm*

\begin{proof}
  Since the views not reachable from~$t$ are the same in~$m$ and~$\hat m$ before the transition,
  and these views remain unchanged, it suffices to show that the memories~$m'$ and~$\hat m'$ resulting from the transition from~$m$ and~$\hat m$, respectively, are $t$-equivalent.
  The transition from check mode $\checkmode$ is derived by the \sem{check} rules,
  so we can apply \cref{lem:simvisit}
  if we can show that the views in $\checkmode$ have no \dec{Effect} decision.
  Note that~$m$'s $t$-semi-stability and $t$-coherence have already been shown in \cref{lem:semistab,lem:coherence}.

  It also follows from \cref{lem:simvisit} that we have $\omega' = \omega''$ if component bodies do not print.

  We show that only views reachable from the root can contain an \dec{Effect} decision, and such views never contain an \dec{Effect} in check mode $\checkmode$ or event loop mode $\eloopmode$.
  \begin{description}
    \item[{\normalfont\Rule{StepInit}}] Every view added by \sem{init} is reachable from the root.
    \item[{\normalfont\Rule{StepEffect}}] The \sem{commitEffs} rules remove all \dec{Effect} decisions from reachable views.
    \item[{\normalfont\Rule{StepCheck}}] Only paths that are passed as arguments to \sem{check} or returned by \sem{reconcile} may have their views contain an \dec{Effect} decision,
    and those views are all reachable after \sem{check} or \sem{reconcile}.
    If any view receives an \dec{Effect} decision, the \sem{check} rule produces the rendered mode $\rendermode$, to which the state transitions.
    \item[{\normalfont\Rule{StepEvent}}] The \field{children} and \field{decision} fields are not modified during the \phase{Normal} phase. \qedhere
  \end{description}
\end{proof}

\chapter{Comparison of Reactive UI Frameworks}\label{app:chap:frameworks}
The categorization of reactive GUI web frameworks based on three properties---(a)~whether they re-read component specifications for re-rendering, (b)~how state updates are processed (queued or immediate), and (c)~which reactivity primitives they employ---are given in \cref{tab:frameworks}.

\begin{table}[tbp]
  \centering
  \caption{Comparison of reactive UI frameworks.}\label{tab:frameworks}
  \smaller[1]
  \begin{tabular}{llll}
    \toprule
    Framework & Re-read Spec. & State Updates & Reactivity Primitives \\
    \midrule
    React & Yes & Queued & Hooks \\
    Preact & Yes & Queued & Hooks \\
    Dioxus & Yes & Immediate & Signals \\
    Solid & No & Immediate & Signals \\
    Leptos & No & Immediate & Signals \\
    Angular & No & Immediate & Signals \\
    Vue & No & Immediate & Signals \\
    Svelte w/ runes & No & Immediate & Signals \\
    Svelte w/o runes & No & Immediate & Compiler-assisted \\
    SwiftUI & Yes & Immediate & Compiler-assisted \\
    \bottomrule
  \end{tabular}
\end{table}

\DefineShortVerb{\+}
We compare various reactive UI frameworks to
\begin{enumerate}
  \item check if state updates are queued or immediate, and
  \item check if component logic gets re-evaluated every render.
\end{enumerate}
\section{React}
+counter+ gets printed every render.
+count+ updates are queued.
\begin{reactcode}
import { useState, useEffect } from 'react';
function Counter() {
  console.log('counter');
  const [count, setCount] = useState(0);
  useEffect(() => {
    if (count >= 3) {
      setCount(0);
    }
  }, [count]);
  function h() {
    console.log(count);
    setCount(count + 1);
    console.log(count);
  }
  return (
    <button onClick={h}>
      {count}
    </button>
  );
}
\end{reactcode}


\section{Preact}
+counter+ gets printed every render.
+count+ updates are queued.
\begin{reactcode}
import { useState, useEffect } from 'preact/hooks';
function Counter() {
  console.log('counter');
  const [count, setCount] = useState(0);
  useEffect(() => {
    if (count >= 3) {
      setCount(0);
    }
  }, [count]);
  function h() {
    console.log(count);
    setCount(count + 1);
    console.log(count);
  }
  return (
    <button onClick={h}>
      {count}
    </button>
  );
}
\end{reactcode}


\section{Dioxus}
+counter+ gets printed every render.
+count+ updates are immediate.
\begin{reactcode}
use dioxus::prelude::*;
fn log(msg: &str) { ... }
#[component]
fn Counter() -> Element {
  log("Counter");
  let mut count = use_signal(|| 0);
  use_effect(move || {
    if count() >= 3 {
      count.set(0);
    }
  });
  rsx! {
    button {
      onclick: move |_| {
        log(count.to_string().as_str());
        count += 1;
        log(count.to_string().as_str());
      },
      "{count}"
    }
  }
}
\end{reactcode}


\section{Solid}
+counter+ gets printed only once.
+count+ updates are immediate.
\begin{reactcode}
import { createSignal, createEffect } from 'solid-js';
function Counter() {
  console.log('counter');
  const [count, setCount] = createSignal(0);
  createEffect(() => {
    if (count() >= 3) {
      setCount(0);
    }
  });
  function h() {
    console.log(count());
    setCount(count() + 1);
    console.log(count());
  }
  return (
    <button onClick={h}>
      {count()}
    </button>
  );
}
\end{reactcode}


\section{Leptos}
+counter+ gets printed only once.
+count+ updates are immediate.
\begin{reactcode}
use leptos::leptos_dom::logging::*;
use leptos::prelude::*;
#[component]
pub fn Counter() -> impl IntoView {
  console_log("Counter");
  let (value, set_value) = signal(0);
  Effect::new(move |_| {
    if value.get() >= 3 {
      set_value.set(0);
    }
  });
  view! {
    <button on:click=move |_| {
      console_log(&value.get().to_string());
      set_value.update(|value| *value += 1);
      console_log(&value.get().to_string());
    }>{value}</button>
  }
}
\end{reactcode}


\section{Angular}
+counter+ gets printed only once.
+count+ updates are immediate.
\begin{reactcode}
import {Component, signal, effect} from '@angular/core';
@Component({
  selector: 'app-root',
  template: `
    <button (click)="h()">
      {{ count() }}
    </button>
  `,
})
export class CounterComponent {
  count = signal(0);
  constructor() {
    console.log('counter');
    effect(() => {
      if (this.count() >= 3) {
        this.count.set(0);
      }
    });
  }
  h() {
    console.log(this.count());
    this.count.update(val => val + 1);
    console.log(this.count());
  }
}
\end{reactcode}


\section{Vue}
+counter+ gets printed only once.
+count+ updates are immediate.
\begin{reactcode}
<script setup>
import { ref, watch } from 'vue';
console.log('counter');
const count = ref(0);
watch(count, (newValue) => {
  if (newValue >= 3) {
    count.value = 0;
  }
});
function h() {
  console.log(count.value);
  count.value += 1;
  console.log(count.value);
}
</script>
<template>
  <button @click="h">
    {{ count }}
  </button>
</template>
\end{reactcode}


\section{Svelte}
\subsection{Svelte with Runes}
+counter+ gets printed only once.
+count+ updates are immediate.
\begin{reactcode}
<script>
  console.log('counter');
  let count = $state(0);
  $effect(() => {
    if (count >= 3) {
      count = 0;
    }
  });
  function h() {
    console.log(count);
    count += 1;
    console.log(count);
  }
</script>
<button onclick={h}>
  {count}
</button>
\end{reactcode}

\subsection{Svelte without Runes}
+counter+ gets printed only once.
+count+ updates are immediate.
\begin{reactcode}
<svelte:options runes={false} />
<script>
  console.log('counter');
  let count = 0;
  $: if (count >= 3) {
    count = 0;
  }
  function h() {
    console.log(count);
    count += 1;
    console.log(count);
  }
</script>
<button on:click={h}>
  {count}
</button>
\end{reactcode}


\subsection{SwiftUI}
+counter+ gets printed every render.
+count+ updates are immediate.
\begin{reactcode}
struct Counter: View {
  @State private var count = 0
  var body: some View {
    print("counter")
    return Button("\(count)") {
      print(count)
      count += 1
      print(count)
    }.onChange(of: count) { oldValue, newValue in
      if newValue >= 3 {
        count = 0
      }
    }
  }
}
\end{reactcode}

\UndefineShortVerb{\+}
